\documentclass[11pt]{article}

\usepackage[T1]{fontenc}
\usepackage{lmodern}
\usepackage[margin=0.82in]{geometry}
\usepackage{microtype}
\usepackage{amsmath,amssymb}
\usepackage{array,booktabs,tabularx}
\usepackage{float}
\usepackage{xcolor}
\usepackage{hyperref}

\definecolor{navy}{HTML}{17324D}
\definecolor{blue}{HTML}{24588D}
\definecolor{teal}{HTML}{16716A}
\definecolor{soft}{HTML}{F1F6F9}
\definecolor{line}{HTML}{D5E1E8}
\hypersetup{
  colorlinks=true,
  linkcolor=blue,
  urlcolor=blue,
  pdftitle={Terminus: Selected Architecture and Research Progress for Offline Indic AI},
  pdfauthor={Compact Intelligence}
}
\setlength{\parindent}{0pt}
\setlength{\parskip}{0.45em}
\renewcommand{\arraystretch}{1.16}
\newcommand{\flowbox}[2]{\fcolorbox{line}{soft}{\parbox[c][1.28cm][c]{#1}{\centering\small\color{navy}#2}}}

\begin{document}
\begin{center}
{\LARGE\bfseries\color{navy} Terminus: Selected Architecture and Research Progress\\[0.14em]
for Offline Indic AI}\par
\vspace{0.4em}
{\large Compact Intelligence}\par
\vspace{0.15em}
{\color{blue} Project progress attachment \quad | \quad October 2026}
\end{center}
\vspace{0.2em}
{\color{teal}\rule{\linewidth}{1.1pt}}

\begin{abstract}
Terminus is a selected, implemented decoder-only transformer architecture
for a compact language model intended for offline Indic-language assistance.
It has about one and a half billion trainable parameters and combines a
four-local-to-one-global attention cadence, grouped-query attention,
Peri-LN normalization, and a learned cross-layer value residual. The training
design starts with dense weights and FP8 matrix computation, then adapts
eligible projection weights through ternary quantization-aware training.
The program targets \textbf{one trillion valid training-token presentations}
across both stages. Its initial Indic-language set is Hindi, Bengali, Tamil,
Telugu, Kannada, Marathi, and Gujarati. This attachment specifies the
architecture first, then records research progress, intended applications,
and the measurements required before release claims.
\end{abstract}

\section{Project definition: selected model architecture}

\subsection{Geometry and end-to-end information path}

Terminus is a causal, decoder-only transformer with twenty-five residual
blocks. Each block contains grouped-query self-attention followed by a
SwiGLU feed-forward network. The frozen tokenizer maps text to a vocabulary
of 65,536 pieces. The input embedding and output classifier share one
weight table; a final RMSNorm precedes next-token prediction. The starting
training context is 4,096 tokens.

\begin{table}[H]
\centering\small
\caption{Selected architecture and dense-stage implementation.}
\begin{tabularx}{\linewidth}{@{}>{\bfseries\raggedright\arraybackslash}p{0.38\linewidth}X@{}}
\toprule
Component & Selected specification \\
\midrule
Model geometry & 25 decoder blocks; hidden width 2,304; SwiGLU width 5,760 \\
Attention geometry & 18 query heads; six key/value heads; head dimension 128; uniform 3:1 query-to-KV ratio \\
Attention range & 20 causal sliding-window layers at 1,024 tokens; five full-context layers at 4,096 tokens \\
Token interface & Frozen 65,536-piece vocabulary; tied embedding and classifier \\
Normalization & Input RMSNorm before attention and FFN; output-side Peri-LN; learned-gain Q/K RMSNorm \\
Cross-layer path & Learned mixture of first-layer and current-layer values in later attention blocks \\
Dense compute & FP32 projection masters; block-scaled FP8 E4M3 linear operands; BF16 residual stream \\
Kernels and loss & FlashAttention-3; DeepGEMM; packed QKV and MLP projections; streaming cross-entropy \\
Trainable size & About one and a half billion parameters; exact selected count: 1,500,456,496 \\
\bottomrule
\end{tabularx}
\end{table}

Of the exact count, 1,349,222,400 parameters are projection weights eligible
for the later ternary path; 150,994,944 belong to the tied
embedding/readout, and 239,152 are other learned parameters. The tied
table remains outside the selected ternary projection set. These counts
separate the compactable part of the model from the weights that currently
retain higher precision.

\begin{center}
\flowbox{0.17\linewidth}{\textbf{Token IDs}\\tied embedding}
\hspace{0.35em}$\longrightarrow$\hspace{0.35em}
\flowbox{0.38\linewidth}{\textbf{Twenty-five transformer blocks}\\
four local + one global, repeated five times\\
Peri-LN + learned value residual}
\hspace{0.35em}$\longrightarrow$\hspace{0.35em}
\flowbox{0.17\linewidth}{\textbf{Next token}\\final norm + tied head}
\end{center}
\vspace{-0.3em}
\begin{center}\small\color{blue}Figure 1. Implemented model path; application-level retrieval and user interface sit outside the model.\end{center}

\subsection{Grouped-query attention and receptive field}

Every layer uses 18 query heads and six key/value heads, so each KV head
serves three query heads. Each head has dimension 128. Query and key
vectors receive learned-gain RMS normalization and rotary position encoding
before attention. The attention operation is
\[
  \operatorname{Attn}_{\ell}(Q,K,V)
  =\operatorname{softmax}\!\left(
     \frac{QK^{\mathsf T}}{\sqrt{128}}+M_{\ell}
   \right)V ,
\]
where the causal mask $M_{\ell}$ either admits the full starting context or
restricts each position to the latest 1,024 tokens. Blocks five, ten,
fifteen, twenty, and twenty-five use full-context attention; all other
blocks use the local mask. The same head geometry applies to both kinds of
block.

This cadence gives local layers a bounded KV-cache window while each fifth
layer can connect distant parts of a passage. Cached inference therefore
allocates a 1,024-token capacity to each local layer and a 4,096-token
capacity to each global layer at the specified context length. Grouped-query
attention reduces KV storage compared with using a separate KV head for
every query head. Extending context beyond 4,096 tokens would require a
separate model and evaluation decision.

\subsection{Peri-LN residual block and SwiGLU feed-forward path}

Let $R_{\gamma}(z)=\gamma\odot z/
\sqrt{\operatorname{mean}(z^2)+\epsilon}$ be RMS normalization with a learned
gain. For block $\ell$, $x_\ell$ is the incoming residual stream, and the
selected computation is
\begin{align}
  u_\ell &= R_{\gamma^{\mathrm{a,in}}_\ell}(x_\ell),\\
  a_\ell &= x_\ell+
    R_{\gamma^{\mathrm{a,out}}_\ell}
    \!\left(W^O_\ell\,\operatorname{Attn}_{\ell}(u_\ell)\right),\\
  h_\ell &= R_{\gamma^{\mathrm{f,in}}_\ell}(a_\ell),\\
  x_{\ell+1} &= a_\ell+
    R_{\gamma^{\mathrm{f,out}}_\ell}
    \!\left(W^D_\ell[
      \operatorname{SiLU}(W^G_\ell h_\ell)\odot W^U_\ell h_\ell
    ]\right).
\end{align}
The two output-side RMSNorm operations sit immediately before residual
addition and have gains distinct from the input normalizers. This is the
Peri-LN component. The FFN uses the gated SwiGLU product shown above,
with intermediate width 5,760. The selected path does not use the older
SubLN branch. Output projection weights use a reduced residual-branch
initialization scale; this is an initialization setting, not an inference
multiplier.

\subsection{Learned cross-layer value residual}

The first attention block's value representation is reused in later blocks.
If $v_0$ is the value vector from the first block and $v_\ell$ is the
current block's value vector, then for later blocks
\[
  \widetilde v_\ell
  =\lambda_{\ell,0}v_0+\lambda_{\ell,1}v_\ell,
  \qquad \ell\geq 1.
\]
The two coefficients are independent learned scalars, each initialized to
one half; they are not constrained to sum to one. Mixing occurs before
attention masking and aggregation. Training preserves a gradient path to
the first block's value projection. In cached inference, the corresponding
first-block values for new positions must remain available alongside the
ordinary per-layer KV caches. This path changes information flow without
changing the local/global schedule.

\section{Training architecture and compression path}

\subsection{Dense FP8 starting stage}

The selected first stage retains FP32 master projection weights. DeepGEMM
uses block-scaled FP8 E4M3 weight operands and FP8 activations for projection
matrix operations, with a full FP8 projection backward path. The residual
stream is BF16. FlashAttention-3 handles attention where supported; packed
QKV and gate/up matrices, fused pointwise operations, and streaming
cross-entropy reduce kernel and memory overhead. The tied output head is
not treated as a ternary projection matrix. The objective is next-token
cross-entropy on tokenized text.

FP8 here describes the arithmetic used for dense projections. It does not
mean the model's trainable master weights or optimizer state are stored
only in FP8. The selected production recipe freezes these model and
optimizer semantics before a run-dependent batch, data manifest, and
schedule are generated.

\subsection{Ternary quantization-aware continuation}

The implemented continuation path converts eligible projection weights,
while keeping trainable master weights. For each 128-by-128 weight tile
$B$, the scale and forward weight are
\[
  s_B=\max\!\left(\frac{1}{|B|}\sum_{i,j\in B}|W_{ij}|,10^{-6}\right),
  \qquad
  Q_B(W)=s_B\,\operatorname{clip}
     \!\left(\operatorname{round}(W_B/s_B),-1,1\right).
\]
The resulting weight codes are minus one, zero, and plus one, multiplied
by their tile scale. One scale covers 16,384 weights. FP8 activations
remain in the linear compute path, and a straight-through weight-gradient
rule updates the FP32 masters. The scale is recomputed from weights, not
trained as an independent parameter. The tied vocabulary table,
normalization gains, and value-mixture coefficients are outside the
ternary projection set.

The dense checkpoint, its immediate ternary conversion, and its trained
ternary continuation must be compared on the same held-out material.
The QAT learning rates, phase duration, data mixture, and optimizer-state
policy require their own reviewed run configuration; dense-stage settings
are not automatically QAT settings. A compact projection representation
also does not establish loaded device RAM: the tied table, scales,
activations, KV cache, and runtime workspace still contribute.

\subsection{Optimizer and distributed execution}

The selected dense optimizer applies spectral Muon to each \emph{logical}
hidden matrix, even when QKV or gate/up projections share a physical packed
parameter. If $W$ has shape $(d_{\rm out},d_{\rm in})$ and $U_t$ is the
Newton--Schulz direction from its momentum-adjusted gradient, the update is
\[
  W_{t+1}=(1-\eta_t\lambda)W_t
    -\eta_t\sqrt{d_{\rm out}/d_{\rm in}}\,U_t .
\]
The aspect-ratio factor applies to each individual Q, K, V, gate, up,
output, and down matrix; weight decay uses the base learning rate.
AdamW handles the tied table and smaller learned gain and scalar groups.

\begin{table}[H]
\centering\small
\caption{Selected dense-stage optimizer starting settings.}
\begin{tabularx}{\linewidth}{@{}>{\bfseries\raggedright\arraybackslash}p{0.70\linewidth}X@{}}
\toprule
Setting & Value \\
\midrule
Spectral Muon base learning rate & 0.017 \\
Tied embedding/readout AdamW learning rate & 0.003 \\
Ordinary normalization, Peri-LN, final norm, value-mix AdamW rate & 0.0003 \\
Learned Q/K norm AdamW rate & 0.0001 \\
Hidden-matrix weight decay; tied-table decay & 0.02; zero \\
Gradient clip; schedule & 1.0; warmup--stable--decay with cosine decay \\
\bottomrule
\end{tabularx}
\end{table}

Replicated data parallelism is the initial multi-GPU strategy: every
rank retains the complete model and optimizer state, accumulates local
microsteps, and reduces FP32 gradients before clipping and updating.
Global batch and token horizons are bound to an external production
manifest. Full-cluster speed, stability, and the dense/QAT cost split
remain measurements to make on the intended hardware, not properties
established by the architecture.

\section{Research progress and training scope}

The project has moved from comparing sequence-model alternatives to a
selected transformer definition and optimized training implementation.
Smaller dense and ternary proxy runs informed the conversion risk;
smaller-width experiments informed optimizer starting settings. The
selected architecture, tokenizer integration, FP8 kernels, distributed
trainer, checkpointing, and a separate QAT path are implemented. Short
complete-update measurements exist on one H100-class GPU; sustained
full-cluster throughput has not been measured.

An earlier proxy trained on about ten billion dense tokens and continued
for about two billion ternary-training tokens. Its immediate conversion
raised held-out prediction loss; continuation recovered much of the loss,
but its final measured result remained below its dense parent on the
tracked held-out mixture. Because the training mixture also changed
between stages, this observation motivates controlled transfer tests
without predicting final quality for the selected larger model.

The program target is \textbf{one trillion valid training-token
presentations across dense learning and ternary adaptation}. A
presentation counts one use of a text token by the trainer. Padding and
held-out evaluation tokens do not count; repeated exposure and distinct
source material are reported separately. The production manifests record
the actual language/source mixture and stage allocation.

The seven initial Indic-language targets are \textbf{Hindi, Bengali,
Tamil, Telugu, Kannada, Marathi, and Gujarati}. English, code, mathematics,
and relevant code-mixed material remain in broader training and
evaluation. Data preparation includes source and license review, script
checks, deduplication, sensitive-data screening, and train/evaluation
separation. Each target language will be evaluated separately.

\section{Applications enabled by the intended offline release}

The selected model is being designed for local-language text assistance
without requiring a live model server for each request. The following
are application designs for an evaluated offline release. They do not
assert that the final trained model or device application has been shipped.
When information can change or carries high stakes, the application
supplies current, trusted local documents and identifies their source.

\begin{table}[H]
\centering\small
\caption{Application designs and the surrounding software they require.}
\begin{tabularx}{\linewidth}{@{}>{\bfseries\raggedright\arraybackslash}p{0.19\linewidth}>{\raggedright\arraybackslash}X>{\raggedright\arraybackslash}X@{}}
\toprule
Application & Example assistance & Application-side requirement \\
\midrule
Public-service forms & Explain a locally stored official form or scheme notice in the user's language and identify the relevant fields. & A versioned official copy, document date and source, and a route to the responsible authority for eligibility decisions. \\
Education & Explain a lesson passage in simpler language and generate practice questions from supplied material. & Source text, age-appropriate controls, and learner or teacher review. \\
Small-business records & Interpret a stock list or invoice draft and explain spreadsheet fields in a target language. & Local files, deterministic calculation of totals, and verified tax or pricing rules where relevant. \\
Field and community work & Summarize an approved local manual as a short checklist in the worker's language. & Current manuals and human review for safety-critical decisions. \\
Personal documents & Search and summarize a user's own stored documents, with pointers to relevant passages. & A local document index and explicit file-access boundaries. \\
\bottomrule
\end{tabularx}
\end{table}

For example, a Hindi-speaking user opens a public-service form that is
already on the device. Local retrieval selects the form's instruction
paragraph; Terminus explains that paragraph in Hindi and cites it back to
the user. If eligibility is not settled by the document, the interface
directs the user to the official channel. Speech input, automatic
submission, and live access to government databases would be separate
software projects with separate evaluation.

\section{Evidence required for completion}

Model evaluation will compare the dense checkpoint, immediate ternary
conversion, and final QAT checkpoint on matched held-out data. It will
report each target language separately, including native script and
relevant code mixing, plus English, code, mathematics, factuality, and
safety-sensitive tasks. The production record will report valid token
exposure, loss and stability traces, checkpoint recovery, and complete
distributed-update speed. Offline readiness requires an exported model
and target-device measurements of file size, peak RAM, response speed,
energy use, and task accuracy.

The current result is a concrete selected architecture and training
implementation, supported by smaller-scale research and one-GPU
engineering checks. The one-trillion-token production run, dense-to-QAT
transfer at the selected scale, full-cluster throughput, multilingual
release quality, and on-device performance remain to be measured.

\end{document}
